\documentclass[11pt,a4paper]{article}

\usepackage[margin=2.5cm]{geometry}
\usepackage[utf8]{inputenc}
\usepackage[T1]{fontenc}
\usepackage{hyperref}
\usepackage{listings}
\usepackage{xcolor}
\usepackage{booktabs}
\usepackage{enumitem}
\usepackage{parskip}
\usepackage{titling}

\definecolor{codebg}{rgb}{0.95,0.95,0.95}
\definecolor{codecomment}{rgb}{0.35,0.45,0.35}
\definecolor{codekeyword}{rgb}{0.05,0.15,0.55}

\lstdefinestyle{cpp}{
    language=C++,
    backgroundcolor=\color{codebg},
    basicstyle=\ttfamily\small,
    keywordstyle=\color{codekeyword}\bfseries,
    commentstyle=\color{codecomment}\itshape,
    stringstyle=\color{purple},
    numbers=left,
    numberstyle=\tiny\color{gray},
    breaklines=true,
    frame=single,
    tabsize=4,
    showstringspaces=false
}

\hypersetup{
    colorlinks=true,
    linkcolor=blue,
    urlcolor=blue,
    citecolor=blue
}

\title{\textbf{UE5\_Spout2\_DX12} \\ \large Technical Changes vs.\ the Original Repository}
\author{Internal engineering note}
\date{2026-07-07}

\begin{document}
\maketitle

\begin{abstract}
\noindent
This document describes, at a technical level, every modification applied
to the \texttt{Spout2\_DX12} plugin on top of the base version tracked in
\texttt{origin/main} (commit \texttt{37e714e}, ``Readme update and hot fix
for viewport sender'' lineage). The changes were developed and validated in
a separate working copy of the plugin embedded in an Unreal Engine 5.8
project, then ported back into this repository. They fall into four
categories: a GPU-hang fix for headless (\texttt{-RenderOffScreen})
sessions, a Shipping-build shutdown crash fix, a receiver-side
frame-synchronization fix, and a linker fix for the third-party Spout
libraries. A command-line build script was also added.
\end{abstract}

\tableofcontents

\section{Scope and Baseline}

\begin{itemize}
    \item \textbf{Baseline}: commit \texttt{37e714e} (source-identical to
        \texttt{baf099f}, which only touches \texttt{README.md}).
    \item \textbf{Files touched}: 7 source files, 476 insertions / 9
        deletions.
    \begin{itemize}[nosep]
        \item \texttt{Source/Spout2\_DX12/Private/Spout2\_DX12.cpp}
        \item \texttt{Source/Spout2\_DX12/Private/SpoutReceiverComponent.cpp}
        \item \texttt{Source/Spout2\_DX12/Private/SpoutSenderComponent.cpp}
        \item \texttt{Source/Spout2\_DX12/Public/Spout2\_DX12.h}
        \item \texttt{Source/Spout2\_DX12/Public/SpoutReceiverComponent.h}
        \item \texttt{Source/Spout2\_DX12/Public/SpoutSenderComponent.h}
        \item \texttt{Source/Spout2\_DX12/Spout2\_DX12.Build.cs}
    \end{itemize}
    \item \textbf{Binaries}: not carried over as-is; recompiled against
        this repository's copy of the source via the new \texttt{Build.bat}
        (Section~\ref{sec:build}), verified to produce a successful
        Editor / Game Development / Game Shipping build for Win64 against
        UE~5.8.
\end{itemize}

\section{Fix 1 --- Deferred Bridge Start (Headless GPU Hang)}
\label{sec:deferred-start}

\subsection{Symptom}
In windowless sessions launched with \texttt{-RenderOffScreen}, the process
intermittently hung with \texttt{DXGI\_ERROR\_DEVICE\_HUNG} on the D3D12
graphics queue, with GPU crash-dump breadcrumbs consistently stuck at
\emph{Frame~0}. The failure rate observed during validation was
approximately 3 out of 4 runs before the fix, 0 out of 4 after.

\subsection{Root Cause}
Both \texttt{USpoutSenderComponent} and \texttt{USpoutReceiverComponent}
open a \texttt{D3D11On12} interop bridge (via \texttt{spoutDX12}) to share
textures with other Spout-aware processes. This bridge submits interop
work on the same D3D12 queue the engine uses for rendering. In a normal
windowed session, the swapchain's \texttt{Present} call implicitly paces
that queue every frame. In a windowless session there is no swapchain, so
nothing paces the queue during the engine's startup phase --- which is also
when the Pipeline State Object (PSO) precompilation flood happens (map
load, first-frame shader compiles). Opening the interop bridge and
submitting work while that flood is in flight was racing the queue and
deadlocking it.

\subsection{Fix Design}
The actual bridge initialization in both components is now held behind two
conditions, checked at the top of \texttt{StartBroadcastConfigured()} (sender)
and \texttt{StartReceiving()} (receiver):

\begin{enumerate}
    \item \textbf{Minimum grace period}: world time must have reached
        \texttt{StartDelaySeconds} (a new \texttt{UPROPERTY}, default
        \texttt{2.0f}; \texttt{0.0f} disables the whole mechanism and
        restores pre-fix immediate-start behavior).
    \item \textbf{PSO precaching idle}: both
        \texttt{PipelineStateCache::NumActivePrecacheRequests()} and
        \texttt{FShaderPipelineCache::NumPrecompilesRemaining()} must be
        zero.
\end{enumerate}

A hard cap, \texttt{MaxStartDelaySeconds} (new \texttt{UPROPERTY}, default
\texttt{15.0f}; \texttt{0.0f} = no cap), forces the start even if PSOs are
still compiling once world time exceeds it --- this is a safety valve so a
scene with a pathologically long shader-compile tail cannot hold the
sender/receiver back indefinitely.

While deferred, a one-shot \texttt{FTimerHandle}
(\texttt{StartDelayTimerHandle}) re-invokes the same entry point after a
retry delay: the remaining grace-period time if still inside it, otherwise
a fixed \texttt{0.5\,s} poll interval. The mechanism only engages in an
actual game world (\texttt{World->IsGameWorld()}); the editor/PIE path is
unaffected because a window's \texttt{Present} already paces the queue
there.

\begin{lstlisting}[style=cpp,caption={Deferred-start gate (sender), abridged}]
UWorld* World = GetWorld();
if (World && World->IsGameWorld() && StartDelaySeconds > 0.0f)
{
    const float WorldSeconds = static_cast<float>(World->GetTimeSeconds());
    const bool bMinTimePassed = WorldSeconds >= StartDelaySeconds;

    const int32 PendingPSOs =
        static_cast<int32>(PipelineStateCache::NumActivePrecacheRequests()) +
        static_cast<int32>(FShaderPipelineCache::NumPrecompilesRemaining());

    const bool bForcedStart =
        MaxStartDelaySeconds > 0.0f && WorldSeconds >= MaxStartDelaySeconds;

    if (!bForcedStart && (!bMinTimePassed || PendingPSOs > 0))
    {
        // stash config, arm a retry timer, bail out without
        // opening the D3D11On12 bridge
        return;
    }
    // otherwise: proceed with the real bridge initialization
}
\end{lstlisting}

\subsection{Lifecycle Cleanup}
\texttt{StopBroadcastInternal()} / \texttt{StopReceivingInternal()} now
explicitly clear \texttt{StartDelayTimerHandle} via
\texttt{World->GetTimerManager().ClearTimer(...)}. Without this, stopping
the sender/receiver while a start was deferred would leave an armed timer
that silently restarted the broadcast/receive path up to \texttt{0.5\,s}
later, independent of the explicit stop. A companion flag,
\texttt{bLoggedStartDeferral}, throttles the ``holding start'' log line to
once per deferral cycle instead of once per \texttt{0.5\,s} retry, and is
reset on stop so a future deferral logs again.

\subsection{Files and Symbols}
\begin{center}
\begin{tabular}{@{}ll@{}}
\toprule
File & New symbols \\
\midrule
\texttt{SpoutSenderComponent.h} & \texttt{StartDelaySeconds},
    \texttt{MaxStartDelaySeconds}, \texttt{StartDelayTimerHandle},
    \texttt{bLoggedStartDeferral} \\
\texttt{SpoutSenderComponent.cpp} & gate in
    \texttt{StartBroadcastConfigured()}; cleanup in
    \texttt{StopBroadcastInternal()} \\
\texttt{SpoutReceiverComponent.h} & same four symbols (independent
    per-component state) \\
\texttt{SpoutReceiverComponent.cpp} & gate in \texttt{StartReceiving()};
    cleanup in \texttt{StopReceivingInternal()} \\
\bottomrule
\end{tabular}
\end{center}

\section{Fix 2 --- Shipping Shutdown Crash}
\label{sec:shutdown-crash}

\subsection{Symptom}
Access violation observed during process teardown in monolithic Shipping
builds.

\subsection{Root Cause}
\texttt{FSpout2\_DX12Module::ShutdownModule()} used to call
\texttt{FPlatformProcess::FreeDllHandle(MyDllHandle)} to explicitly unload
\texttt{SpoutDX12.dll}. In monolithic Shipping builds, static destructors
running \emph{after} module shutdown can still call into that DLL; once
it has been unloaded, that call lands on unmapped memory.

\subsection{Fix}
The explicit \texttt{FreeDllHandle} call was removed;
\texttt{MyDllHandle} is simply reset to \texttt{nullptr} as a bookkeeping
value. The DLL is left mapped for the remaining lifetime of the process ---
harmless, since the OS reclaims all module handles unconditionally at
process exit. This trades an unnecessary explicit unload for guaranteed
safety against late static-destructor calls.

\section{Fix 3 --- Graceful Stop Channel for Headless Sessions}
\label{sec:stop-channel}

\subsection{Problem}
A windowless process has no \texttt{WM\_CLOSE} to receive, so there is no
clean way to ask it to exit from the outside. Terminating it with
\texttt{taskkill /F} skips Unreal's orderly shutdown path
(\texttt{EndPlay} $\rightarrow$ \texttt{StopBroadcastInternal} $\rightarrow$
Spout sender release), leaving a ``ghost'' sender registered in Spout's
shared memory. The next launch's sender then gets auto-renamed
(\texttt{"Sender Component\_1"}, \texttt{"\_2"}, \ldots) by Spout's naming
collision handling, breaking any receiver (LiveBoard, OBS, \ldots) hard-wired
to the original fixed name.

\subsection{Fix: File-Based Command Channel}
\texttt{StartupModule()} now computes a per-project stop-file path:

\begin{lstlisting}[style=cpp]
StopFilePath = FPaths::Combine(
    FString(FPlatformProcess::UserTempDir()),
    FString::Printf(TEXT("%s_SpoutStop"), FApp::GetProjectName()));
\end{lstlisting}

i.e.\ \texttt{\%TEMP\%\textbackslash{}<ProjectName>\_SpoutStop}. Any stale
file from a previous run is deleted immediately on startup (otherwise a new
session would see it and self-terminate roughly one second in). A
\texttt{FTSTicker} callback (\texttt{CheckStopFile}), registered on the
engine's core ticker at a 1-second period, polls for that file's existence.
The core ticker runs even with no viewport, so this works in headless mode.

When the file is detected:
\begin{enumerate}
    \item It is deleted immediately (one-shot semantics --- otherwise it
        would also kill the \emph{next} session right after startup).
    \item \texttt{RequestEngineExit(...)} is called --- Unreal's
        \emph{orderly} exit path, as opposed to a hard kill: it lets the
        current frame finish, tears the world down, and lets
        \texttt{EndPlay}/\texttt{StopBroadcastInternal} release the Spout
        sender from shared memory before the process actually exits
        ($\sim$1.5\,s in practice).
\end{enumerate}

Two guards prevent unwanted behavior:
\begin{itemize}
    \item \texttt{GIsEditor} --- the channel is inert inside the editor.
        A file dropped in \texttt{\%TEMP\%} must never be able to close an
        open Editor session; the channel is only meaningful for
        packaged/game sessions.
    \item \texttt{IsEngineExitRequested()} --- avoids requesting exit twice
        if it is already in progress.
\end{itemize}

The ticker handle (\texttt{StopFileTickerHandle}) is explicitly removed in
\texttt{ShutdownModule()} \emph{before} the module object is destroyed,
since the delegate was bound with \texttt{CreateRaw} (no lifetime
tracking) --- a tick firing after that point would call into freed memory.

An external tool (e.g.\ a \texttt{.bat} script, not part of this
repository) can request a clean shutdown of a running headless instance
simply by doing:
\begin{lstlisting}[language=bash]
type nul > "%TEMP%\<ProjectName>_SpoutStop"
\end{lstlisting}

\section{Fix 4 --- Receiver Frame Synchronization}
\label{sec:frame-sync}

\subsection{Problem}
The receiver copied the sender's shared D3D11 texture on every tick with
no synchronization against the sender's write cadence, risking tearing or
reading a texture mid-write when sender and receiver frame rates diverge.

\subsection{Fix}
\texttt{USpoutReceiverComponent} now tracks a small piece of Spout
frame-sync state:

\begin{lstlisting}[style=cpp]
bool bSyncInitialized = false;
char SpoutSyncName[256] = {};
\end{lstlisting}

In \texttt{ReceiveOnce()}, once the active sender's name is known, the
mutex/frame-count objects are (re)initialized whenever the sender identity
changes (first connection, or a different sender name than last time):

\begin{lstlisting}[style=cpp]
if (!bSyncInitialized || FCStringAnsi::Strcmp(SpoutSyncName, SenderName) != 0)
{
    if (bSyncInitialized)
    {
        SpoutDX12->frame.CloseAccessMutex();
        SpoutDX12->frame.DisableFrameCount();
    }
    SpoutDX12->frame.CreateAccessMutex(SenderName);
    SpoutDX12->frame.EnableFrameCount(SenderName);
    FCStringAnsi::Strncpy(SpoutSyncName, SenderName, 256);
    bSyncInitialized = true;
}
\end{lstlisting}

Two checks then guard the actual texture copy:
\begin{itemize}
    \item \textbf{New-frame check}: \texttt{SpoutDX12->frame.GetNewFrame()}
        is a non-blocking read of the sender's frame semaphore; if the
        sender has not produced a new frame since the last read, the
        receiver skips the copy this tick. The very first frame
        (\texttt{bConnected == false}) always bypasses this check so the
        initial image is never dropped.
    \item \textbf{Access mutex}: \texttt{SpoutDX12->frame.CheckAccess()} is
        acquired (non-blocking) immediately before
        \texttt{Ctx11->CopyResource(...)}, and released with
        \texttt{AllowAccess()} immediately after. If the sender is
        mid-write, \texttt{CheckAccess()} returns \texttt{false}, the copy
        is skipped for that tick, and a missed-frame counter
        (\texttt{StatMissedFramesAccum}) is incremented for diagnostics.
\end{itemize}

Sync state is torn down symmetrically in both
\texttt{StopReceivingInternal()} and \texttt{ReleaseSpoutDevices()}
(\texttt{CloseAccessMutex()} / \texttt{DisableFrameCount()}, then
\texttt{bSyncInitialized = false} and \texttt{SpoutSyncName[0] = '\textbackslash 0'}), so a
subsequent \texttt{StartReceiving()} reinitializes cleanly rather than
inheriting stale mutex/semaphore handles.

\section{Fix 5 --- Third-Party Linker Configuration}
\label{sec:linker}

\subsection{Symptom}
Crash at runtime tied to \texttt{spoutFrameCount} symbol resolution.

\subsection{Root Cause}
\texttt{SpoutDX12.dll} statically embeds the \texttt{spoutDX} /
\texttt{spoutFrameCount} implementation. The module's
\texttt{Spout2\_DX12.Build.cs} previously linked \emph{both}
\texttt{Spout.lib} and \texttt{SpoutDX12.lib}, and delay-loaded \emph{both}
\texttt{Spout.dll} and \texttt{SpoutDX12.dll}. Linking \texttt{Spout.lib}
caused the linker to resolve \texttt{spoutFrameCount} symbols against
\texttt{Spout.dll} instead of the statically-embedded copy inside
\texttt{SpoutDX12.dll}. Since \texttt{Spout.dll} is never actually loaded
at runtime by \texttt{SpoutDX12.dll}, any call through those symbols
crashed.

\subsection{Fix}
\begin{lstlisting}[basicstyle=\ttfamily\small,backgroundcolor=\color{codebg},frame=single,breaklines=true]
- PublicAdditionalLibraries.Add(Path.Combine(LibDir, "Spout.lib"));
  PublicAdditionalLibraries.Add(Path.Combine(LibDir, "SpoutDX12.lib"));

- PublicDelayLoadDLLs.AddRange(new[] { "Spout.dll", "SpoutDX12.dll" });
+ PublicDelayLoadDLLs.Add("SpoutDX12.dll");
\end{lstlisting}

Only \texttt{SpoutDX12.lib} is linked, and only \texttt{SpoutDX12.dll} is
delay-loaded; \texttt{Spout.dll} is still staged next to the packaged
executable as a runtime dependency (unrelated third-party consumers may
still need the standalone Spout runtime), but the plugin's own code no
longer resolves symbols against it.

\section{Addition --- Command-Line Build Script}
\label{sec:build}

A new \texttt{Build.bat} wraps
\texttt{RunUAT.bat BuildPlugin}, the official Epic tooling for compiling a
standalone plugin without an open Editor session (this is the same
mechanism used to produce Marketplace plugin binaries). It takes two
optional positional arguments, \texttt{EngineDir} (default
\texttt{C:\textbackslash Program Files\textbackslash Epic
Games\textbackslash UE\_5.8}) and \texttt{OutputDir} (default
\texttt{.\textbackslash Build\textbackslash Spout2\_DX12}), and invokes:

\begin{lstlisting}[language=bash]
RunUAT.bat BuildPlugin -Plugin=<uplugin path> -Package=<OutputDir> \
    -TargetPlatforms=Win64 -Rocket
\end{lstlisting}

This compiles three Win64 configurations --- Editor, Game Development, Game
Shipping --- and stages a self-contained, redistributable plugin folder
(\texttt{.uplugin}, \texttt{Binaries/}, \texttt{Source/},
\texttt{Resources/}) at \texttt{OutputDir}. The script was validated
end-to-end against UE~5.8: all three configurations compiled and linked
successfully (\texttt{BUILD SUCCESSFUL}), total wall time $\approx$
2\,minutes on the reference machine. Full usage and installation
instructions (per-project vs.\ engine-wide) are documented separately in
\texttt{BUILD\_INSTRUCTIONS.md}.

\section{Summary Table}

\begin{center}
\begin{tabular}{@{}p{0.28\linewidth}p{0.32\linewidth}p{0.32\linewidth}@{}}
\toprule
Change & Failure mode addressed & Mechanism \\
\midrule
Deferred bridge start & GPU hang in
    \texttt{-RenderOffScreen} sessions & world-time + PSO-idle gate,
    capped retry \\
No explicit DLL unload & Shipping shutdown access violation &
    leave \texttt{SpoutDX12.dll} mapped until process exit \\
File-based stop channel & Ghost sender after \texttt{taskkill /F} &
    polled stop-file $\rightarrow$ \texttt{RequestEngineExit} \\
Receiver frame sync & Tearing / mid-write reads & access mutex +
    frame-count semaphore around \texttt{CopyResource} \\
Linker fix & Crash resolving
    \texttt{spoutFrameCount} & drop \texttt{Spout.lib} /
    \texttt{Spout.dll} delay-load \\
\texttt{Build.bat} & No command-line build path & wraps
    \texttt{RunUAT.bat BuildPlugin} \\
\bottomrule
\end{tabular}
\end{center}

\end{document}
